% Registro de decisiones: una fila por decisión importante (la más nueva arriba).
\clearpage
\section{Registro de decisiones}

\begin{center}
\small
\begin{tabularx}{\linewidth}{>{\raggedright\arraybackslash}p{.15\linewidth}>{\raggedright\arraybackslash}X>{\raggedright\arraybackslash}X}
\toprule
\textbf{Fecha} & \textbf{Decisión} & \textbf{Por qué} \\
\midrule
10-nov-2026 & Usar los conceptos del curso Proyectos como guía para formalizar Franui, sin que sea una tarea. & Las herramientas (MAP, doble diamante, validación) sirven para construir la startup y postular a fondos. \\
10-nov-2026 & Llevar esta bitácora en LaTeX, dentro del repositorio del proyecto. & Queda guardada con su historial de cambios y se puede compartir como PDF. \\
31-oct-2026 & Estética minimalista para la app y la web, con la misma paleta. & Referencias de marcas premium; más elegante y fácil de leer. \\
31-oct-2026 & Quitar del banco las expresiones universales (\emph{oh, oh dear, wow, hooray}). & No se aprenden: se entienden en cualquier idioma. \\
14-oct-2026 & Mudar la app a \url{https://polykat-franui.github.io}. & Una dirección con el nombre del proyecto, sin perder el avance de nadie. \\
14-oct-2026 & Toda traducción generada por IA queda como borrador hasta revisarla. & La calidad del contenido es la base de la confianza. \\
\bottomrule
\end{tabularx}
\end{center}
